%% kapa: exam=1 round=24 year=2013 period=2 form=A subject=appraisal_law subject_ko=부동산관계법규 seq=26 number=66
%% source: Q-Net 감정평가사 기출문제 (cst00309) artlSeq=5073802 file=extracted/1008/5073802/2013년 제24회 감정평가사 1차시험 기출문제/2교시(A형).hwp
%% conversion: kordoc 4.24.51 -> tools/gen1.py (auto); review=proofread-sonnet; flags=none
\begin{question}{66}
국유재산법령상 국유재산의 대장과 보고에 관한 설명으로 옳은 것은?
\begin{choices}
  \item 국방부장관은 그가 관리하는 항공기가 멸실된 경우에 지체없이 그 사실을 총괄청과 감사원에 보고하여야 한다.
  \item 총괄청으로부터 일반재산의 관리ㆍ처분에 관한 사무를 위탁받은 투자매매업자가 해당 사무와 관련하여 등기소의 장에게 등기사항증명서의 교부를 청구하려면 수수료를 납부하여야 한다.
  \item 총괄청은 매년 국유재산관리운용보고서와 검사보고서를 감사원에 제출하여야 한다.
  \item 총괄청은 중앙관서별로 국유재산에 관한 총괄부를 갖추어 두어 그 상황을 명백히 하여야 하기 때문에 그 총괄부를 전산자료로 대신할 수 없다.
  \item 중앙관서의 장은 국유재산의 구분과 종류에 따라 그 소관에 속하는 국유재산의 대장ㆍ등기사항증명서와 도면을 갖추어 두어야 하며, 국유재산의 대장은 전산자료로 대신할 수 있다.
\end{choices}
\end{question}
\answer{5}
